\begin{abstract}
\noindent
Large-language-model (LLM) agents that read and modify code pay two
recurring taxes: a \emph{read tax} when they consume entire source files
to answer questions whose answer lies in the file's structure rather than
its bodies, and a \emph{write tax} when they emit free-form diffs that can
silently corrupt working programs. Both taxes scale with file size; both
can be cut by routing through the abstract syntax tree (AST). This paper
introduces and evaluates the design of an AST-native code surface for the
multi-agent framework \blackrim: a single Markdown outline emitter with
three entry points (a CLI, a runtime hook that auto-prepends outlines to
\texttt{Read} tool results, and an MCP retrieval surface), a
plan/execute pair contract for symbolic surgery, and a polyglot
compile-gate that rejects refactors whose dry-run output fails the
language's native syntax check.

We formalise the design as a constrained reduction problem: minimise the
expected token cost of an agent's interaction with a file while
preserving task-success rate $\success$ and ruling out false-positive
edits (those that compile or parse without error but alter behaviour
unintendedly). The reduction has three measurable components: outline
compression ratio $\ratio(\file) = |\outline(\file)|/|\file|$, outline
adoption rate (the fraction of large-file \texttt{Read}s preceded by a
\texttt{gt outline} call), and refactor false-negative rate (the fraction
of symbolic edits whose compile-gate verdict matches a ground-truth
review). The first is measured today; the second and third are the
subject of pre-registered empirical questions \texttt{OQ-AST-1..OQ-AST-7}
in \cref{app:open-questions}.

Implemented across Go, Python, JavaScript, and TypeScript via per-language
parser backends (\texttt{go/ast}, embedded CPython \texttt{ast},
\texttt{acorn}, \texttt{ts.createSourceFile}), our compression backends
achieve measured token savings of \textbf{82.0\% (Go)}, \textbf{84.1\%
(Python)}, \textbf{83.0\% (JavaScript)}, and \textbf{75.2\%
(TypeScript)} on a realistic benchmark fixture, with graceful
fallback-to-passthrough on parse error or missing host runtime. The
outline emitter targets a $\sim 300$-token output regardless of file
size, yielding an estimated $\sim 100\times$ token reduction on a
2700-line file relative to a full \texttt{Read}; the end-to-end
reduction on real agent traces is the central pending empirical claim.

We discuss limitations --- in particular the tree-sitter grammar
versioning surface, the prompt-injection threat posed by feeding
verbatim file content (sanitised) into the outline lens, and the
fidelity--velocity trade-off introduced by a YAML pattern-DSL for
authoring new transform intents --- and pre-register the open
questions that gate promotion from outline-as-tool to
outline-as-default.

\textbf{Keywords:} abstract syntax tree, code outline, context window,
multi-agent systems, LLM tooling, tree-sitter, code refactoring, language
server protocol, prompt compression.
\end{abstract}
